% sec_identification.tex — V10: honest pre-trends discussion

\subsection{Additional Identification Evidence}
\label{sec:additional_id}

\paragraph{FL entry event study.}Within-PBU event studies around first FL entry
(Figure~\ref{fig:fl_entry_es} in Appendix~\ref{app:robustness}) show
positive pre-event coefficients at $t = -2$ and $t = -3$, indicating
that PBUs receiving FL entry \emph{already had higher prices}. This
pattern admits two interpretations:

\begin{enumerate}
  \item \emph{Strategic market selection}: cartels select into
    competitive, high-value markets---the pre-trend reflects this
    selection---and then deploy cover bidders to sustain elevated prices.
    Under this interpretation, the pre-trend is consistent with the
    cover-bidding hypothesis and the OLS fixed effects absorb the
    \emph{level} of prices in each market.
  \item \emph{Selection bias}: unobserved market characteristics that
    attract both FL firms and higher prices generate a spurious
    association. Under this interpretation, the item, year, and PBU
    fixed effects may not fully absorb time-varying within-market
    confounders, and the OLS overstates the FL--price association.
\end{enumerate}

\noindent We cannot definitively distinguish between these
interpretations. Two pieces of evidence favor the first: the
reverse-causality test (Section~\ref{sec:mechanisms}, M3) shows
that the elasticity of FL entry with respect to lagged prices is
economically negligible ($\approx 0.004$), bounding the magnitude of
price-chasing selection. The sensitivity analysis ($RV = 17.5\%$)
implies that a confounder driving the pre-trend would need to explain
a substantial share of residual variation in both FL presence and log
prices to fully account for the association. Neither test eliminates
selection bias; we regard the pre-trends as an unresolved ambiguity
that reinforces treating the estimates as conditional associations
rather than causal effects.

\paragraph{Minimum-bidder constraint variation.}
In convite, when $n_{\text{genuine}} < 3$ the cartel \emph{must}
deploy cover bidders (corner solution). When $n \geq 3$, cover bidding
is voluntary. Interacting FL with a constraint-binding indicator within
the convite subsample yields $\hat{\beta}_{\text{FL}}(n \geq 3) =
0.076$ (SE~$= 0.021$, $p < 0.001$)---a 7.6\% premium from
\emph{voluntary} cover bidding---and $\hat{\beta}_{\text{FL} \times
(n<3)} = -0.160$ ($p < 0.001$). The price effect concentrates in
voluntary deployments, not constraint-driven ones.

\paragraph{CADE enforcement shock.}
The Callaway--Sant'Anna DiD using CADE conviction timing shows that FL
presence declines significantly post-conviction (ATT~$= -0.275$,
$p < 0.001$), consistent with cover bidders exiting after cartel
prosecution. The price ATT is positive but imprecise (0.145,
SE~$= 0.110$).
